\subsection{面向 Decode 稳态性能的离线 Priming 机制}

\subsubsection{设计动机}

在前述分阶段调度基础上，进一步实验发现：当运行时不再在主流程中显式执行算子提交参数预热，而是直接进入真实请求时，Decode 吞吐会对首次执行路径表现出较强敏感性。具体而言，在同样的正式请求配置下，若模型在进入正式测量前已经经历过一次较大线程池上的 Prefill/Decode 执行，则后续正式 Decode 的吞吐更高，且波动更小；反之，若直接在冷态下进入正式请求，则 Decode 吞吐较低，并且在部分配置下会出现首次执行阶段异常不稳定的问题。

这一现象说明，影响 Decode 稳态性能的因素并不完全来自正式请求本身的算子计算，还包括运行时首次建立的线程池状态、算子提交路径、页映射与访存相关状态等。换言之，正式请求观察到的吞吐差异，部分来源于首次运行态初始化成本是否已经被提前支付。

\subsubsection{基本思路}

基于上述观察，本文在不破坏正式请求主流程的前提下，引入一类一次性的隐藏预热操作，即离线 Priming。其核心思想是：在进入正式请求前，先使用一组可控参数执行一次不计入结果统计的隐藏 Prefill/Decode，用于显式触发目标运行路径上的运行时初始化；随后立即恢复正式阶段的线程数、绑核配置与调度参数，再开始真实请求测量。

该机制并不改变正式请求的输入输出，也不改变正式请求的线程数与绑核配置，仅将原本发生在首次正式请求中的初始化成本提前到一次隐藏执行中完成。因此，它本质上属于运行时稳态固化机制，而不是新的调度算法。

\subsubsection{机制实现}

本文实现的 Priming 机制包含以下四个特征。

\begin{enumerate}
  \item \textbf{一次性执行。} Priming 只在正式测量前执行一次，不参与统计。
  \item \textbf{参数可控。} 隐藏预热可独立设置 prompt 长度、decode token 数、Prefill 线程数与 Decode 线程数。
  \item \textbf{执行后恢复正式计划。} Priming 完成后立即执行状态重置，并恢复正式请求对应的阶段线程数、绑核配置和调度参数。
  \item \textbf{默认关闭。} 若不显式开启该机制，原有正式请求主流程保持不变。
\end{enumerate}

从工程实现角度看，该机制可抽象为式~(\ref{eq:decode-prime-flow}) 所示流程：

\begin{equation}
\label{eq:decode-prime-flow}
\text{Load} \rightarrow \text{Phase Plan Setup} \rightarrow \text{Hidden Prime}
\rightarrow \text{Reset \& Restore} \rightarrow \text{Formal Request}
\end{equation}

其中，Hidden Prime 阶段仅负责提前触发目标执行路径，不计入正式实验结果；Formal Request 阶段仍采用论文前文确定的正式参数。

\subsubsection{为何能够同时改善 Decode 与 Prefill}

实验进一步表明，该机制不仅会提高后续 Decode 吞吐，在部分长 prompt 场景下也会提高后续正式 Prefill 吞吐。但这一收益并不意味着正式 Prefill 算子本身被单独优化。相反，消融实验显示，只要隐藏预热能够以目标大线程池路径完整执行一次，即使其 prompt 长度仅为 1 token，后续正式 Prefill 仍会变快。这说明收益主要来自运行态首次初始化被提前触发，而不是大 prompt 数据本身被“跑热”。

因此，本文将该现象解释为：离线 Priming 预先建立了与大线程池执行路径相关的运行时状态，包括线程池工作状态、算子提交路径以及页映射/地址转换相关状态，正式 Prefill 与 Decode 均可复用这些已初始化的状态，从而减少首次正式请求中的额外开销。

\subsubsection{适用边界}

需要强调的是，离线 Priming 的价值主要体现在\textbf{会话级}或\textbf{模型常驻}场景，而不适合作为\textbf{每个请求前置执行}的固定步骤。原因在于，Priming 本身需要额外执行一次隐藏 Prefill，若将其放在每次请求前执行，则会增加首请求时延，并可能抵消短回复场景下的吞吐收益。

因此，本文建议将该机制部署为：模型加载完成后的单次后台预热，或会话建立后的单次冷启动预热。这样，预热成本可以被后续多个正式请求分摊，而正式请求本身无需再承担首次执行路径的冷态初始化成本。

\subsubsection{与前文调度策略的关系}

离线 Priming 与前文提出的分阶段线程数配置、真实性能比初始划分以及 Prefill work stealing 机制并不冲突。前者解决的是\textbf{首次执行路径的运行态初始化问题}，后者解决的是\textbf{稳态执行阶段的线程组织与负载均衡问题}。因此，本文将其视为对阶段感知调度框架的工程补充，而不是对原有调度算法的替代。

